\section{Proofs for the critical length spectrum}\label{app:critical}

This appendix proves Theorem~\ref{thm:critical-cluster}, following the
outline in Subsection~\ref{sec:critical-cluster}. Throughout,
$n/\sqrt g\to c\in(0,\infty)$, and $P_s$, $K_s$, $\mu_c$ and $\nu$ are as
in \eqref{eq:cluster-intensities}.

\subsection{Ordinary cuts and cusp allocations}

Short multicurves are organized by their cut graphs: the vertices are the
components of the complement, labeled by their genera, and the edges are
the curves. Lemma~\ref{lem:critical-cuts} bounds the expected number of
multicurves with a given cut graph. It shows that, in the limit, only
nonseparating curves and boundaries of two-cusp pairs of pants
contribute; the later subsections compute these contributions jointly.

Fix a connected cut graph with $k\ge1$ edges and $Q$ vertices;
loops and multiple edges are allowed. At vertex $i$ let $a_i$ be the
genus and $p_i\ge1$ its degree. Then
\[
 \sum_i a_i=g-k+Q-1,\qquad \sum_i p_i=2k.
\]
Partition the vertices into $u$ of type $(0,1)$, $v$ of type $(0,2)$,
and $w$ with $h_i=2a_i-2+p_i\ge1$. Put $h_0=2g-2$. For large
$g$, $w\ge1$, and the latter complexities sum to $H=h_0+u$.
A component of type $(0,1)$ needs at least two cusps, and one of
type $(0,2)$ at least one.

\begin{lemma}[Critical cut estimate]\label{lem:critical-cuts}
Fix $k$ and a bound on the cut lengths. At $n/\sqrt g\to c>0$,
the expected number of cut multicurves of the specified graph and
vertex partition, summed over all admissible genera and cusp allocations,
is $O(g^{1-w-v/2})$. If a specified $(0,1)$ component has at least
three cusps, the bound gains a factor $g^{-1/2}$.
\end{lemma}

\begin{proof}
Write $h_n=2g-2+n$ for the complexity of the cusped surface. The volume bounds in \cite[Lemmas~18--19]{SW} give
\[
 V_{a,b}\le C\frac{(4\pi^2)^{d-1}\Gamma(d)}{\sqrt d},
 \quad d=2a-2+b\ge1,\qquad
 V_{g,n}\ge c_0\frac{(4\pi^2)^{h_n-1}\Gamma(h_n)}{\sqrt{h_n}}.
\]
The first bound includes $d=1$ after enlarging $C$ for $(0,3)$ and
$(1,1)$; the second uses $n=O(\sqrt g)$. If vertex $i$ receives
$q_i$ cusps, put $d_i=h_i+q_i$. Then $\sum_i d_i=h_n$ and
$\sqrt{h_n}/\prod_i\sqrt{d_i}\le\sqrt Q$. Mirzakhani integration,
the bounded-boundary estimate \cite[Lemma~3.2]{Mir}, and the
multinomial allocation of cusp labels therefore bound the expected
cut count $E_{g,n}$ by
\[
 E_{g,n}\le\frac{C_{k,c}\,n!}{\Gamma(h_n)}
       \sum_{\substack{\text{admissible genera}\\q_1+\cdots+q_Q=n}}
       \prod_i\frac{\Gamma(h_i+q_i)}{q_i!}.
\]
Here and below the constants may also depend on the length bound.
The powers of $4\pi^2$ leave only the fixed factor $(4\pi^2)^{1-Q}$.

The cusp sums at the two unstable vertex types are bounded using
\begin{equation}
 \phi_{0,1}(z)=\sum_{q\ge2}\frac{z^q}{q(q-1)}=\frac{z^2}{2}+O(z^3),
 \qquad \phi_{0,2}(z)=\sum_{q\ge1}\frac{z^q}{q}=z+O(z^2).
 \label{eq:cluster-unstable}
\end{equation}
At a stable vertex their generating function is
$\Gamma(h)(1-z)^{-h}$. Moreover
\[
 \sum_{\substack{h_1+\cdots+h_w=H\\h_i\ge1}}
       \prod_i\Gamma(h_i)\le C_w\Gamma(H-w+1).
\]
For $w=2$ this follows by dividing by $\Gamma(H-1)$ and summing
$\binom{H-2}{h_1-1}^{-1}$; iteration gives the general bound.
Dropping genus parity constraints only increases the sum.
Let $a_m=[z^m]\phi_{0,1}(z)^u\phi_{0,2}(z)^v$. Expanding $(1-z)^{-H}$ gives
\[
 E_{g,n}\le C_{k,c}\frac{\Gamma(H-w+1)}{\Gamma(H)}
       \sum_{m\le n}a_m(n)_m\frac{\Gamma(h_n+u-m)}{\Gamma(h_n)}.
\]
Uniformly for $m\le n=O(\sqrt{h_n})$, the finite-product formula for
the Gamma ratio gives
$\Gamma(h_n+u-m)/\Gamma(h_n)\le C_{k,c}h_n^{u-m}$: the logarithmic
error is $O((m+1)^2/h_n)$. Since $(n)_m\le n^m$, positivity yields
\begin{equation}
 E_{g,n}\le C_{k,c}h_n^{u-w+1}\phi_{0,1}(n/h_n)^u\phi_{0,2}(n/h_n)^v
          =O(g^{1-w-v/2}).
 \label{eq:critical-graph-bound}
\end{equation}
Requiring three cusps at a specified $(0,1)$ vertex replaces one
$\phi_{0,1}(z)$ by $\phi_{0,1}(z)-z^2/2$ and gains $n/h_n=O(g^{-1/2})$.
\end{proof}

Only $w=1$, $v=0$ can survive: the graph then has one large-genus
vertex, loops at that vertex, and two-cusp pants attached as leaves.
All other graphs, and all leaves with more than two cusps, vanish.

\subsection{Excluding complex short subsurfaces}\label{subsec:critical-exclusion}

We need expectation bounds, not only a bound on the probability of a
bad configuration, so we give the multiplicity argument explicitly.

\begin{lemma}[Bounded primitive count]\label{lem:primitive-count}
For fixed $K,L$, a finite-area hyperbolic surface $S$ with geodesic
boundary and $0<-\chi(S)\le K$ has at most $C(K,L)$ primitive
unoriented geodesics of length at most $L$, uniformly in its metric.
\end{lemma}

\begin{proof}
Double $S$ along its boundaries. The double has area at
most $4\pi K$. Fix a Margulis constant $\varepsilon>0$. A primitive
geodesic wholly in a thin component is its core; there are boundedly
many such cores. Every other geodesic meets the thick part. A maximal
$\varepsilon/4$-separated net there has bounded cardinality by area
packing. At each net point, orbit points of a lift are uniformly
separated, so hyperbolic ball packing bounds the number of based
loops of length at most $L+\varepsilon$. Connecting a short geodesic
to a net point proves the assertion. Primitivity is preserved upon
inclusion in the double or in an ambient surface: a root would have
the same geodesic image and would already be a closed curve in $S$.
\end{proof}

\begin{lemma}[Complex short configurations]\label{lem:critical-complex}
Fix $r$ and $L$. As $g\to\infty$ with $n/\sqrt g\to c\in(0,\infty)$,
\begin{equation}
 \mathbb E[\text{intersecting ordered simple $r$-tuples of lengths
 at most $L$}]=O_{r,L,c}(g^{-1/2}),
 \label{eq:critical-intersections}
\end{equation}
\begin{equation}
 \mathbb E[\text{nonsimple primitive curves of length at most $L$
 not filling two-cusp pants}]=O_{L,c}(g^{-1/2}),
 \label{eq:critical-bad-curves}
\end{equation}
and the expected number of simple separating curves of length at most $L$
that are not waists of two-cusp pairs of pants tends to zero.
\end{lemma}

\begin{proof}
For fixed $r,L$, a tuple of $r$ distinct simple geodesics of length
at most $L$ has bounded total intersection number. We use the local
segment argument in the proof of \cite[Proposition~4.5]{MP}; it
applies equally in the presence of cusps. Thicken each nonannular
component of their union, fill complementary disks and peripheral
once-punctured disks, and tighten the essential boundaries. The
resulting pieces have bounded topology and total boundary length at
most $2rL$. Keep one ambient geodesic for each parallel boundary class.
Empty annuli between parallel copies collapse; their two incidences
are the two sides of a single edge, possibly a loop. An annulus with
cusps does not collapse and is a $(0,2)$ component counted among the
$v$ vertices.
Isolated tuple curves are retained as additional cut edges, unless
they are already one of the tightened boundaries. Minimal intersection
ensures that these boundaries are disjoint from the original tuple.

There are finitely many resulting graph shapes and markings at fixed
$r,L$. For a fixed cut multicurve, Lemma~\ref{lem:primitive-count}
bounds the choices of original curves in each marked piece. Hence the
tuple count is bounded pointwise by a constant times a finite sum of
marked cut-multicurve counts. An intersecting simple tuple cannot fill
two-cusp pants. Its marked component either has $h_i\ge1$, forcing
a second stable vertex to carry the unbounded genus, or has type
$(0,2)$, or has type $(0,1)$ with at least three cusps. The bounds in
Lemma~\ref{lem:critical-cuts} give \eqref{eq:critical-intersections}.

For a single nonsimple primitive curve of length at most $L$, use
the filling construction of \cite[Definition~4.10 and Lemma~4.11]{Hide}.
The filled piece has boundary length at most $2L$ and area at most
$4L$, so its topology is bounded. The same recording retains formal
boundary sides that tighten to the same ambient curve as sides of a
cut edge. Lemma~\ref{lem:primitive-count} controls the number of curves
filling each piece. The preceding classification therefore gives
\eqref{eq:critical-bad-curves}.
A filled piece cannot be the whole ambient surface for large $g$,
since its topology is bounded. Finally, the one-edge case of
Lemma~\ref{lem:critical-cuts} shows that simple separating curves
of bounded length other than two-cusp pants waists vanish in expectation.
\end{proof}

\subsection{Joint nonseparating and pants counts}

Let $\mathcal N_g$ count simple nonseparating lengths and let
$\mathcal Q_g$ count two-cusp pants waist lengths.

\begin{lemma}[Joint Poisson limit]\label{lem:critical-skeleton}
As $g\to\infty$ with $n/\sqrt g\to c\in(0,\infty)$,
\begin{equation}
 (\mathcal N_g,\mathcal Q_g)\Longrightarrow
 (\mathcal N,\mathcal Q_c)
 \label{eq:critical-skeleton}
\end{equation}
jointly in the vague topology on $[0,\infty)$, where $\mathcal N$ and
$\mathcal Q_c$ are independent Poisson point processes of intensities $\nu$
and $\mu_c$.
\end{lemma}

\begin{proof}
In the mixed
factorial moments of $(\mathcal N_g,\mathcal Q_g)$, \eqref{eq:critical-intersections} removes intersecting
tuples. Lemma~\ref{lem:critical-cuts} removes every disjoint cut type
except the graph with a single large component, $k$ nonseparating
loops and $r$ two-cusp leaves. For prescribed bounded intervals for
the $k+r$ lengths, its exact contribution is
\begin{equation}
 \frac{(n)_{2r}}{2^{r+k}V_{g,n}}
 \int V_{g-k,n-r+2k}(\mathbf s,\mathbf t,\mathbf t,0^{n-2r})
        \prod_{j=1}^r s_j\,\mathrm ds_j
        \prod_{i=1}^k t_i\,\mathrm dt_i.
 \label{eq:critical-mixed-moment}
\end{equation}
Here $(n)_{2r}$ is the falling factorial. The factor $2^{-r}$ counts
unordered pairs of cusp labels at the ordered pants leaves, as in
\cite[Lemma~28]{SW}; $2^{-k}$ exchanges the sides of each
nonseparating cut, as in \cite[Proposition~4.2]{MP}. There is no
factor $r!$ or $k!$, since both tuples are ordered.

The volume ratio is a product of $r$ cusp ratios and $k$ genus ratios.
The cusp ratios satisfy $V_{g,j+1}/V_{g,j}\sim8\pi^2g$ for $j\le n$ by
\cite[Lemma~5.1(ii)]{MZ}, since $n=o(g)$.
For the genus ratios, \cite[Lemma~5.1(iii)]{MZ} gives only an upper bound,
and the two-sided limit at $n\asymp\sqrt g$ is the zero-density case
$\mathcal B(0)=1$ of Theorem~\ref{thm:volume}. Together,
\[
 \frac{V_{g-k,n-r+2k}}{V_{g,n}}\sim(8\pi^2g)^{-r}.
\]
Proposition~\ref{prop:bounded-boundaries}, applied at zero limiting
density, where $\mathfrak b_0(L)=\sinh(L/2)/(L/2)$, supplies, uniformly
on the fixed length box, the boundary factor
$\prod_j S(s_j/2)\prod_i S(t_i/2)^2$, with $S(t)=\sinh(t)/t$; compare
Shen--Wu's bounded-boundary estimate~\cite[Proposition~25]{SW}, which
they attribute essentially to \cite[Proposition~3.1]{MP}
and~\cite[Lemma~20]{NWX}.
Consequently the limiting factors in \eqref{eq:critical-mixed-moment}
are
\[
 \frac12 tS(t/2)^2\,\mathrm dt=\nu(\mathrm dt),\qquad
 \frac{c^2}{16\pi^2}sS(s/2)\,\mathrm ds=\mu_c(\mathrm ds).
\]
All mixed factorial moments on disjoint bounded intervals converge to
those of independent Poisson counts. Higher fixed moments give uniform
integrability, and the multivariate Poisson law is moment determinate.
The first-moment bounds near zero are $O(\varepsilon^2)$. This proves the
lemma.
\end{proof}

\subsection{Proof of Theorem~\ref{thm:critical-cluster}}

\begin{proof}[Completion of the proof of Theorem~\ref{thm:critical-cluster}]
Fix $L<\infty$. By Lemma~\ref{lem:critical-complex}, with probability
tending to one every nonsimple primitive curve of length at most $L$
fills two-cusp pants with waist at most $2L$. By
\eqref{eq:critical-intersections}, the pants waists up to $2L$ are
pairwise disjoint with probability tending to one. The pants interiors
are then disjoint as well. Other simple separating curves up to $L$
are absent with probability tending to one. On this event, the full
primitive counting measure restricted to $[0,L]$ is exactly
\[
 \left.\left(\mathcal N_g+
       \sum_{s\in\mathcal Q_g\cap[0,2L]}K_s\right)\right|_{[0,L]}.
\]
Every interior primitive geodesic of $P_s$ is nonsimple and fills it.
The filling boundary bound implies $s\le2\ell$ for such a geodesic
of length $\ell$. Thus pants with waist greater than $2L$ cannot
contribute, and the limiting sum is locally finite.

For $\varepsilon>0$, the family $P_s$, $\varepsilon\le s\le2L$,
is compact in its fixed-topology moduli space. Lengths of marked
classes vary continuously, and only finitely many such classes have
length at most $L$ somewhere on this compact family. For the latter
assertion, use a finite cover by uniformly bilipschitz markings,
with standard neighborhoods at the two cusps, and the finiteness of
each center metric's length spectrum below a fixed bound. Hence
$s\mapsto\langle f,K_s\rangle$ is continuous away from zero for
continuous compactly supported $f$. Apply \eqref{eq:critical-skeleton}
with waist cutoff $\varepsilon$, choosing $L$ larger than the support
of $f$. The probability of removing any waist is at most
$\mathbb E\mathcal Q_g([0,\varepsilon])$, whose limit is
$\mu_c([0,\varepsilon])=O(\varepsilon^2)$. Letting
$\varepsilon\downarrow0$ proves \eqref{eq:critical-cluster}.
The independent Poisson Laplace functionals give
\eqref{eq:cluster-laplace}.
\end{proof}
